%% ---------------------------------------------------------------------
%%  I Radicali: Semplificazione — la seconda scheda, per provare che la
%%  libreria si riusa.
%%
%%  Guarda il preambolo: due righe e la testata. Nessuna grafica copiata,
%%  nessun colore scelto, nessuna forma disegnata: se qui servisse
%%  ricopiare qualcosa, la libreria non sarebbe pronta.
%% ---------------------------------------------------------------------
\documentclass[10pt]{article}
\usepackage{schede}

\schedaTitolo{I Radicali: Semplificazione}
             {Proprietà invariantiva \& portar fuori}
\schedaPiede{Schema di Matematica — I Radicali (2)}

\begin{document}

%% =====================================================================
\schedaSezione{1}{La proprietà invariantiva}
%% =====================================================================

\riquadro[La regola]{Moltiplicare o dividere per lo stesso numero (diverso
da zero) l'indice della radice e l'esponente del radicando non cambia la
radice: è la stessa radice scritta in un altro modo.}

\vspace{6pt}

\begin{schedaTabella}{Cosa si fa}{Schema visivo (stesso numero $n$)}{Esempio}

  \riga{\voce{Semplificare l'indice}{(dividere per $k$)}}
       {$\radice[n\cdot k]{\formaQuadrato^{k}}
         \;\freccia\;\radice[n]{\formaQuadrato}$}
       {$\radice[6]{\formaQuadrato^{2}}
         =\radice[3]{\formaQuadrato}$}

  \riga{\voce{Portar fuori}{(se il radicando è un prodotto)}}
       {$\radice[n]{\formaQuadrato^{n}\cdot\formaTriangolo}
         \;\freccia\;\formaQuadrato\cdot\radice[n]{\formaTriangolo}$}
       {$\radice[3]{\formaQuadrato^{3}\cdot\formaTriangolo}
         =\formaQuadrato\cdot\radice[3]{\formaTriangolo}$}

  \riga{\voce{Portar dentro}{(il fattore davanti)}}
       {$\formaQuadrato\cdot\radice[n]{\formaTriangolo}
         \;\freccia\;\radice[n]{\formaQuadrato^{n}\cdot\formaTriangolo}$}
       {$\formaCerchio\cdot\radice[2]{\formaTriangolo}
         =\radice[2]{\formaCerchio^{2}\cdot\formaTriangolo}$}

\end{schedaTabella}

\vspace{4pt}

\riquadro{Le tre mosse si leggono al contrario l'una dell'altra: portare
fuori e portare dentro sono la stessa uguaglianza, e semplificare l'indice è
la proprietà invariantiva letta da sinistra a destra.}

\end{document}
